\chapter{Diseño Analítico de la Etapa en Colector Común}
\label{chap:diseno_teorico}

En este capítulo se desarrolla el dimensionamiento en corriente continua de la etapa amplificadora en configuración
colector común, fundamentando la elección de las corrientes y tensiones de polarización y diseñando la red de base
bajo criterios estrictos de estabilidad térmica.

\section{Circuito Esquemático de Partida}

En la figura~\ref{crkt:colector_comun} se presenta el esquema circuital completo del amplificador colector común
(seguidor de emisor) a implementar.

\begin{figure}[!ht]
  \centering
  \begin{circuitikz}[scale=0.95, transform shape]
  % Referencia y fuente de señal
  \draw (0,0) node[ground] {}
    to[sV, l=$v_s$] (0,3)
    to[R, l=$R_s$] (1.8,3)
    to[C, l=$C_1$, -*] (3.5,3) coordinate (base_node);

  % Divisor resistivo de entrada (R1 a masa, R2 a VCC)
  \draw (3.5,6.5) node[vcc] {$V_{CC}$} coordinate (vcc_node)
    to[R, l=$R_2$, -*] (base_node)
    to[R, l=$R_1$, *-] (3.5,0) node[ground] {};

  % Transistor BJT NPN en Colector Común
  \draw (base_node) to[short] (4.2,3) node[npn, anchor=B] (Q1) {};
  \node[right=4pt] at (Q1.right) {$Q_1$ (BC546B)};

  % Colector directamente a VCC (o desacoplado a masa en CA)
  \draw (Q1.C) to[short] ++(0,1.5) coordinate (aux_c) to[short] (vcc_node -| aux_c) node[vcc] {$V_{CC}$};

  % Malla de emisor con resistencia RE
  \draw (Q1.E) to[short, -*] ++(0,-0.6) coordinate (emit_node);
  \draw (emit_node) to[R, l=$R_E$] ++(0,-2.0) node[ground] {};

  % Red de acoplo de salida desde el emisor y carga RL
  \draw (emit_node) to[C, l=$C_2$] ++(3.5,0) coordinate (out_node)
    to[short, -o] ++(1,0) node[above] {$v_L$};
  \draw (out_node) to[R, l=$R_L$, *-] ++(0,-2.0) node[ground] {};
\end{circuitikz}

  \caption{Esquema circuital completo de la etapa adaptadora en configuración colector común.}
  \label{crkt:colector_comun}
\end{figure}

El capacitor de acoplo de entrada $C_1$ aísla el nivel continuo de la base frente al transductor, mientras que $C_2$
bloquea la componente continua de emisor ($V_E$) para entregar únicamente la señal senoidal amplificada a la carga $R_L$.
El colector del transistor $Q_1$ se conecta de manera directa al riel positivo de alimentación $V_{CC}$, constituyendo
una masa dinámica de corriente alterna.

\section{Dimensionamiento del Punto de Operación en Continua}

La etapa se alimenta con una fuente de continua estándar de laboratorio regulada en $V_{CC} = 5.0\V$.

\subsection{Fijación de la Corriente de Colector de Reposo}
Para que la etapa opere en clase A y reproduzca la señal senoidal sin incurrir en recorte dinámico durante los semiciclos
negativos, la corriente estática de emisor debe ser estrictamente superior a la corriente pico demandada por la carga:
\begin{equation}
  I_{CQ} \ge \frac{\hat{v}_L}{R_L} = \frac{\sqrt{2} \cdot 200\mV_{\mathrm{rms}}}{10\kohm} = \frac{282.84\mV}{10000\ohm} \approx 28.28\uA
  \label{eq:Icq_min}
\end{equation}

Si bien una corriente de $30\uA$ evitaría el corte teórico, operar con corrientes tan bajas degradaría el ancho de
banda debido a las capacidades parásitas de juntura, incrementaría la resistencia intrínseca $r_e$ elevando la
impedancia de salida y haría al circuito muy susceptible al ruido térmico. Por tales motivos de diseño práctico,
se adopta:
\begin{equation}
  I_{CQ} = 1.0\mA
\end{equation}

\subsection{Fijación de la Tensión de Emisor y Cálculo de $R_E$}
Para permitir la libre oscilación pico de la tensión sobre el emisor sin aproximarse a masa, se exige que la tensión
continua de reposo satisfaga $V_E \ge \hat{v}_L \approx 0.28\V$. Adoptando un margen de seguridad frente a variaciones
térmicas y tolerancias, se define:
\begin{equation}
  V_E = 1.0\V
\end{equation}

A partir de la ley de Ohm en la rama de emisor:
\begin{equation}
  R_E = \frac{V_E}{I_{EQ}} \approx \frac{V_E}{I_{CQ}} = \frac{1.0\V}{1.0\mA} = 1000\ohm = 1.0\kohm
  \label{eq:RE_calculo}
\end{equation}

La tensión colector-emisor de reposo resultante queda determinada por la malla de salida:
\begin{equation}
  V_{CEQ} = V_{CC} - V_E = 5.0\V - 1.0\V = 4.0\V
  \label{eq:VCEQ_calculo}
\end{equation}
situando al transistor en una región de polarización activa directa óptima, distante tanto de la saturación
($V_{CE,\mathrm{sat}} \approx 0.2\V$) como del corte ($V_{CE} = 5.0\V$).

\section{Diseño de la Red de Polarización de Base}

Para asegurar que el punto de trabajo se mantenga insensible a las dispersiones de fabricación de la ganancia
de corriente ($\beta$) y a las fluctuaciones de temperatura, se aplica el teorema de Thévenin al divisor resistivo
de entrada formado por $R_1$ y $R_2$. En la figura~\ref{crkt:thevenin} se ilustra el circuito equivalente.

\begin{figure}[!ht]
  \centering
  \begin{circuitikz}[scale=0.95, transform shape]
  % Malla Thévenin de base
  \draw (0,0) node[ground] {}
    to[battery1, invert, l=$V_{BB}$] (0,2.5)
    to[R, l=$R_B$, i=$I_{BQ}$] (2.8,2.5)
    to[short] (3.6,2.5) node[npn, anchor=B] (Q1) {};
  \node[right=4pt] at (Q1.right) {$Q_1$};

  % Colector directamente a VCC
  \draw (Q1.C) to[short] (Q1.C |- 0,4.5) node[vcc] {$V_{CC}$};

  % Emisor
  \draw (Q1.E) to[R, l=$R_E$, i=$I_{EQ}$] (Q1.E |- 0,0) node[ground] {};

  % Indicadores de tensión VBE y VCE
  \draw[dashed, color=gray] (3.6,2.4) to[open, v<=$V_{BE}$] (4.3,1.5);
  \draw[dashed, color=gray] (4.8,3.2) to[open, v^<=$V_{CEQ}$] (4.7,1.4);
\end{circuitikz}

  \caption{Circuito equivalente de Thévenin para el análisis de polarización de base en colector común.}
  \label{crkt:thevenin}
\end{figure}

El criterio de diseño estándar de la cátedra para garantizar un divisor de base rígido establece:
\begin{equation}
  R_B \le \frac{\beta}{10} \cdot R_E
  \label{eq:criterio_rb}
\end{equation}

Tomando inicialmente el valor típico de catálogo para el BC548/BC547B ($\beta = 300$):
\begin{equation}
  R_B = \frac{300}{10} \cdot 1000\ohm = 30.0\kohm
\end{equation}

Planteando la ecuación de la malla de entrada de base:
\begin{equation}
  V_{BB} = I_{BQ} \cdot R_B + V_{BE} + I_{EQ} \cdot R_E
  \label{eq:malla_thevenin}
\end{equation}

Sustituyendo $I_{BQ} = I_{CQ}/\beta$ y $R_B = (\beta/10) R_E$:
\begin{equation}
  I_{BQ} \cdot R_B = \frac{I_{CQ}}{\beta} \left(\frac{\beta}{10} R_E\right) = 0.10 \cdot I_{CQ} \cdot R_E
\end{equation}

Aproximando $I_{EQ} \approx I_{CQ}$ con $V_{BE} = 0.70\V$:
\begin{equation}
  V_{BB} = 0.10 \cdot I_{CQ} R_E + I_{CQ} R_E + V_{BE} = 1.10 \cdot I_{CQ} R_E + V_{BE}
\end{equation}
\begin{equation}
  V_{BB} = 1.10 \cdot (1.0\mA \cdot 1000\ohm) + 0.70\V = 1.10\V + 0.70\V = 1.80\V
  \label{eq:VBB_calculada}
\end{equation}

\section{Cálculo de los Resistores del Divisor ($R_1$ y $R_2$)}

Las relaciones de Thévenin para el divisor de tensión con $R_1$ conectado entre base y masa y $R_2$ conectado
entre $V_{CC}$ y base son:
\begin{equation}
  V_{BB} = V_{CC} \cdot \frac{R_1}{R_1 + R_2}, \quad R_B = R_1 \parallel R_2 = \frac{R_1 \cdot R_2}{R_1 + R_2}
\end{equation}

Dividiendo ambas expresiones miembro a miembro:
\begin{equation}
  \frac{V_{BB}}{R_B} = \frac{V_{CC}}{R_2} \implies R_2 = \left(\frac{V_{CC}}{V_{BB}}\right) \cdot R_B
\end{equation}
\begin{equation}
  R_2 = \left(\frac{5.0\V}{1.80\V}\right) \cdot 30.0\kohm \approx 83.33\kohm
  \label{eq:R2_teorico}
\end{equation}

Despejando $R_1$:
\begin{equation}
  R_1 = \frac{R_B}{1 - \frac{V_{BB}}{V_{CC}}} = \frac{30.0\kohm}{1 - 0.36} = \frac{30.0\kohm}{0.64} \approx 46.88\kohm
  \label{eq:R1_teorico}
\end{equation}

\section{Ajuste Analítico para Ganancia Medida ($\beta = 284$)}

En el banco de trabajo, la ganancia estática de corriente del transistor individual seleccionado se midió en
$\beta = 284$. Recalculando con este valor exacto:
\begin{equation}
  R_B = \frac{284}{10} \cdot 1000\ohm = 28.4\kohm
\end{equation}
\begin{equation}
  R_2 = \left(\frac{5.0\V}{1.80\V}\right) \cdot 28.4\kohm \approx 78.89\kohm
\end{equation}
\begin{equation}
  R_1 = \frac{28.4\kohm}{1 - 0.36} \approx 44.38\kohm
\end{equation}

En el Cuadro~\ref{tab:resumen_diseno_teorico} se sintetizan las magnitudes ideales obtenidas para la etapa.

\begin{table}[!ht]
  \centering
  \begin{tabular}{|l|c|c|l|}
    \hline
    \textbf{Parámetro de Diseño} & \textbf{Símbolo} & \textbf{Valor Teórico} & \textbf{Unidad} \\
    \hline
    Tensión de alimentación continua & $V_{CC}$ & $5.00$ & $\V$ \\
    Corriente de colector de reposo & $I_{CQ}$ & $1.00$ & $\mA$ \\
    Tensión continua de emisor & $V_E$ & $1.00$ & $\V$ \\
    Tensión colector-emisor de reposo & $V_{CEQ}$ & $4.00$ & $\V$ \\
    Tensión Thévenin de base & $V_{BB}$ & $1.80$ & $\V$ \\
    Resistencia Thévenin de base & $R_B$ & $28.40$ & $\kohm$ \\
    Resistencia de emisor calculada & $R_E$ & $1.00$ & $\kohm$ \\
    Resistencia inferior del divisor (a masa) & $R_1$ & $44.38$ & $\kohm$ \\
    Resistencia superior del divisor (a $V_{CC}$) & $R_2$ & $78.89$ & $\kohm$ \\
    Corriente de base de reposo & $I_{BQ}$ & $3.52$ & $\uA$ \\
    \hline
  \end{tabular}
  \caption{Resumen del dimensionamiento teórico en corriente continua para la etapa colector común.}
  \label{tab:resumen_diseno_teorico}
\end{table}
